\begin{corollary}[Conditional revised-or-restricted embedding route]
\label{cor:conditional-embedding-route}
Consider a KL localization stage whose underlying map is a physical
embedding \(a\), with source and target localization subgroups
\(H_s,H_t\).  If
\[
 a_*^{-1}(H_t)=H_s,
\]
then the localized inclusion is an embedding on the covered
complementary component.  If this equality holds at every stage through
the complete direct-decoration horizon, then all induced degree-one
universal-cover tile maps are injective; together with the separately
proved component-bijection and tile-gluing criterion, the corresponding
fixed pseudo-piece immersion is an interior holomorphic embedding.
\end{corollary}

\begin{proof}
% The underlying research notes state this corollary without a separate
% proof; each step below is one application of a result proved in this
% paper.
The first assertion is Theorem~\ref{thm:subgroup-reflection-criterion}
applied to the stage's localization covers: the reflection equality
forces the localized lift to be injective, and it is an embedding on
the covered complementary component (an injective holomorphic
immersion is an open injective map, hence a homeomorphism onto its
image).  For the second assertion, stagewise injectivity of the
localized base maps upgrades, via the tile/base equivalence
(Theorem~\ref{thm:tile-base-equivalence}, proved in
Section~\ref{sec:tiles} below), to injectivity of every
induced degree-one universal-cover tile lift.  The final assertion
combines tile injectivity with the fixed-pair component bijection
through the tile-gluing criterion
(Proposition~\ref{prop:tile-glue-equivalence}, also proved in
Section~\ref{sec:tiles} below) to conclude that the
glued pseudo-piece immersion is an interior holomorphic embedding.
No circularity arises from the forward references: nothing in
Sections~\ref{sec:criteria}--\ref{sec:basechange} uses this corollary.
\end{proof}